\documentclass[11pt]{article}
\usepackage[margin=2.1cm]{geometry}
\usepackage{booktabs}
\usepackage{amsmath}
\usepackage{enumitem}
\usepackage[T1]{fontenc}
\usepackage{lmodern}

\setlength{\parindent}{0pt}
\setlength{\parskip}{5pt}
\newcommand{\code}[1]{\texttt{#1}}
\newcommand{\answerbox}[1]{\vspace{2pt}\fbox{\parbox[c][#1][c]{\dimexpr\linewidth-2\fboxsep-2\fboxrule}{\ }}\vspace{4pt}}

\begin{document}

\begin{center}
{\Large\bfseries EE\,362 --- Pre-Lab 2}\\[2pt]
{\large Master--Slave Flip-Flop and the N-Bit Register}\\[6pt]
\begin{tabular}{@{}p{6.5cm}p{4cm}p{4cm}@{}}
Name: \hrulefill & Section: Mon / Tue & Date: \hrulefill
\end{tabular}
\end{center}

\vspace{2pt}\hrule\vspace{4pt}

{\small\textbf{Due at the start of the lab session.} Five points, about twenty
minutes, on paper, before you open Quartus.}

\vspace{4pt}\hrule

\section*{Question 1 --- Paper first}

Draw a timing diagram. The clock runs for two full periods. \code{D} goes high
somewhere in the middle of the first low phase, and goes low again somewhere in
the middle of the second high phase.

Sketch \code{Clk}, \code{D}, the master output $Q_m$, and the final output
\code{Q}. Mark every point where $Q_m$ changes and every point where \code{Q}
changes.

\answerbox{5.2cm}

\textbf{One sentence:} \code{Q} changes at only one kind of instant. Which, and
what physically prevents it changing at any other time?

\answerbox{1.6cm}

\section*{Question 2 --- Interface}

Write the VHDL \code{entity} for \code{RegN}, including its \code{generic}.
Interface only.

\answerbox{4.4cm}

Now write the \emph{sensitivity list} you will use for the process inside
\code{RegN}, and say why each signal is or is not in it. \textbf{Careful:} this
is not the same rule as Lab 1. Consider \code{Clk}, \code{Rst}, \code{Ld} and
\code{D} one at a time.

\answerbox{3.0cm}

\section*{Question 3 --- Test plan}

List the minimum set of vectors that would convince \emph{you} that a design is
edge-triggered and not merely transparent.

\begin{center}
\begin{tabular}{@{}c@{\hskip 8pt}p{5.2cm}@{\hskip 8pt}p{7.4cm}@{}}
\toprule
\# & Stimulus (and \emph{when}, relative to the clock) & What this vector proves \\
\midrule
1 & \rule{5cm}{0.4pt} & \rule{7.2cm}{0.4pt} \\[8pt]
2 & \rule{5cm}{0.4pt} & \rule{7.2cm}{0.4pt} \\[8pt]
3 & \rule{5cm}{0.4pt} & \rule{7.2cm}{0.4pt} \\[8pt]
4 & \rule{5cm}{0.4pt} & \rule{7.2cm}{0.4pt} \\[8pt]
5 & \rule{5cm}{0.4pt} & \rule{7.2cm}{0.4pt} \\[8pt]
6 & \rule{5cm}{0.4pt} & \rule{7.2cm}{0.4pt} \\
\bottomrule
\end{tabular}
\end{center}

\textbf{Then, in one or two sentences:} a transparent latch would pass some of
your vectors. Which of them would it \emph{fail}? That vector is the one doing
the real work --- everything else is bookkeeping.

\answerbox{2.2cm}

\vfill
{\small\emph{Question 3 is the one that matters.} It is the same question as
Section 5 of your lab report.}

\end{document}
